% Generated draft: requires mathematical and source review.
\documentclass[11pt]{article}
\usepackage[margin=25mm]{geometry}
\usepackage{fontspec}
\setmainfont{Latin Modern Roman}
\usepackage{amsmath,amssymb,mathtools,unicode-math}
\setmathfont{Latin Modern Math}
\usepackage{enumitem,needspace,xurl,hyperref}
\setlength{\parindent}{0pt}
\setlength{\parskip}{5pt}
\setlength{\emergencystretch}{3em}
\newcommand{\N}{\mathbb{N}}
\newcommand{\Z}{\mathbb{Z}}
\newcommand{\Q}{\mathbb{Q}}
\newcommand{\R}{\mathbb{R}}
\newcommand{\C}{\mathbb{C}}
\newcommand{\F}{\mathbb{F}}
\newcommand{\A}{\mathbb{A}}
\newcommand{\PP}{\mathbb{P}}
\newcommand{\E}{\mathbb{E}}
\newcommand{\eps}{\varepsilon}
\newcommand{\dd}{\,\mathrm{d}}
\begin{document}
\section*{UnsolvedMath ID: 700113}
\textbf{Source problem code:} catalogue entry
\subsection{Definitions and mathematical statement}
A Lie ring is an abelian group \((L,+,-,0)\) equipped with a biadditive operation \([\ ,\ ]:L\times L\to L\) satisfying
\[
[x,x]=0,\qquad [x,[y,z]]+[y,[z,x]]+[z,[x,y]]=0
\quad\text{for all }x,y,z\in L.
\]
Work in the signature \((+,-,0,[\ ,\ ])\), with no prescribed characteristic or additional scalar field. An identity is a universally quantified equation \(s(\mathbf{x})=t(\mathbf{x})\) between terms in this signature. For a class \(\mathfrak K\) of Lie rings, define its \(I\)-theory by
\[
I(\mathfrak K)=\{s=t: A\models\forall\mathbf{x}\,s(\mathbf{x})=t(\mathbf{x})\text{ for every }A\in\mathfrak K\}.
\]
A quasi-identity is a universally quantified formula of the form
\[
\forall\mathbf{x}\left(\bigwedge_{j=1}^{m}s_j(\mathbf{x})=t_j(\mathbf{x})\ \Longrightarrow\ s(\mathbf{x})=t(\mathbf{x})\right),
\qquad m\geq 0.
\]
The \(Q\)-theory \(Q(\mathfrak K)\) consists of all quasi-identities valid in every member of \(\mathfrak K\). These definitions apply analogously in the signatures of the other algebraic structures mentioned in the source.

A variety of Lie rings is finitely axiomatizable if there is a finite set \(\Sigma\) of identities such that
\[
\mathcal V=\{L:L\text{ is a Lie ring and }L\models\Sigma\}.
\]
Here finite axiomatizability is understood in the usual equational sense for varieties. The defining Lie ring identities themselves form a finite set. The theory \(I(\mathcal V)\) is decidable if an algorithm, given any identity in this signature, always halts and determines whether that identity belongs to \(I(\mathcal V)\), using a fixed effective encoding of terms.

\textbf{Remaining question, part c)(3)(i), proposed by A. I. Mal'cev (1967).} Does there exist a finite set \(\Sigma\) of Lie ring identities for which the variety \(\mathcal V\) defined above has undecidable \(I\)-theory? Equivalently, does there exist a finitely axiomatizable variety of Lie rings whose set of universally valid identities is not recursive?

\textbf{Supplied archival context.} The record also reproduces questions concerning finitely axiomatizable varieties of a) groups, b) semigroups, c)(1) rings, c)(2) associative rings, with subparts (i) concerning undecidable \(I\)-theory and (ii) concerning undecidable \(Q\)-theory, and c)(3)(ii) Lie rings with undecidable \(Q\)-theory. Its trailing phrase asks for non-decidable \(I\)-theory (\(Q\)-theory), but the exact allocation of these two alternatives to a), b), and c)(1) is unclear in the supplied formatting. The archive reports affirmative answers for a), attributed to Yu. G. Kleiman; b), attributed to V. L. Murski\u\i; and c)(1), attributed to V. Yu. Popov. It reports a negative answer for c)(2)(i), attributed to A. Ya. Belov, and affirmative existence answers for c)(2)(ii) and c)(3)(ii), attributed to A. I. Budkin. These are supplied archival assessments, not independently verified conclusions in this draft. The meaning of unrestricted ``rings'' in c)(1), including its signature and axioms, is not specified sufficiently to reconstruct that archived question precisely.

\textbf{Supplied remark.} The source states that it is not difficult to find a recursively axiomatizable variety of semigroups with identity whose \(I\)-theory is non-recursive, and refers also to A. I. Mal'cev's 1966 paper listed below. Here recursively axiomatizable means that an algorithm can enumerate an equational axiom set; it does not require that set to be finite. The source does not specify whether the semigroup identity element is named by a constant in the signature, so this remark is retained without selecting that convention.
\subsection{Short English statement}
The remaining question asks whether finitely many identities can define a variety of Lie rings for which no algorithm decides all identities valid throughout the variety. The supplied record labels the overall problem partially solved and lists archived answers to related questions. Those labels and answers are supplied source assessments; the archival wording leaves some of the related questions ambiguous.
\subsection{Sources}
\begin{itemize}
\item Supplied local UnsolvedMath registry record: permanent ID \(700113\), exact source problem code \texttt{catalogue entry}, title \emph{Kourovka Notebook Problem 2.40}; upstream revision \texttt{d371ee917ac9647fbb06cd6fc529deec6f3dda82}. The record attributes the problem to A. I. Mal'cev, proposed in 1967.
\item \emph{The Kourovka Notebook (Unsolved Problems in Group Theory)}, 2nd issue (1967), Problem 2.40; supplied record reports an update in the 21st issue (2026), arXiv:1401.0300v48. The supplied collection description names E. I. Khukhro and V. D. Mazurov as editors of the Issues 1--21 collection.
\item Supplied status review, checked on \texttt{2026-10-16}: canonical status \texttt{partially\_solved}, evidence field \texttt{published}, with solved parts assigned to the Archive of Solved Problems and remaining parts to the main section. This is supplied dated curation, not independently verified here.
\item Yu. G. Kleiman, \emph{Trans. Moscow Math. Soc.}, 1983, no. 2, pp. 63--110. Supplied archive attribution for the affirmative answer to a).
\item V. L. Murski\u\i, \emph{Math. Notes} \textbf{3} (1968), pp. 423--427. Supplied archive attribution for the affirmative answer to b).
\item V. Yu. Popov, \emph{Math. Notes} \textbf{67} (2000), pp. 495--504. Supplied archive attribution for the affirmative answer to c)(1).
\item A. Ya. Belov, \emph{Izv. Math.} \textbf{74}, no. 1 (2010), pp. 1--126. Supplied archive attribution for the negative answer to c)(2)(i).
\item A. I. Budkin, \emph{Izv. Altai Univ.} \textbf{65}, no. 1 (2010), pp. 15--17 (Russian). Supplied archive attribution for affirmative existence answers to c)(2)(ii) and c)(3)(ii).
\item A. I. Mal'cev, \emph{Mat. Sbornik} \textbf{69}, no. 1 (1966), pp. 3--12 (Russian). Supplied reference accompanying the remark about recursively axiomatizable varieties of semigroups with identity.
\end{itemize}
\end{document}
